\chapter{构思研究问题与制订研究计划}
\chapterauthors{Steven R. Cummings、Alka M. Kanaya}
\sourcepage{43}
研究问题就是研究者希望通过研究消除的不确定性。值得研究的问题很多，即使我们成功回答了一些问题，身边仍然充满其他问题。例如，在老年人中，孤独与功能状态下降及死亡率上升有很强的关联（1）。但这一发现又引出新的问题：提供有组织的居家支持服务，能否减轻老年人的孤独？这种服务能否防止功能下降？同伴支持、家庭支持或小组支持，哪种干预最有效？通过电话或视频远程开展的服务能否奏效？孤独是否也与认知功能下降有关？

寻找研究问题的挑战，在于明确一个重要而尚未解答的问题，并将其转化为可行、合乎伦理且具有有效性的研究计划。本章介绍完成这一任务的策略。

\section{研究问题的来源}
对已有研究经验的研究者，最好的研究问题通常来自她在自己以往的研究及本领域其他科学家的研究中观察到的发现与问题。初学者尚未积累这样的经验。虽然新鲜视角有时能使富于创造力的人为老问题提出新方法，但总体而言，缺乏经验是一种障碍。

一个好的起点是弄清研究问题与研究兴趣的区别。请看下面这个研究问题：
\sourcepage{44}
\begin{studyquote}定期提供基本收入，能否减轻中年成年人的抑郁？\end{studyquote}
提出它的人，研究兴趣可能是健康的社会决定因素、抑郁的预防或管理，或者卫生政策干预。区分研究问题与研究兴趣很重要，因为某个具体问题即使不能转化为可行的研究计划，研究者仍可另提一个问题，继续探究自己感兴趣的领域。

当然，如果连研究兴趣都不确定，就不可能形成研究问题。如果你处于这种状况，并不孤单：许多初学者还没有找到既令自己兴奋、又适合制订可行计划的主题。可以先想一想，阅读时哪些研究曾引起你的兴趣。你也许记得某位患者的治疗似乎不充分或不恰当：如果换一种做法，她的结局是否会改善？也可以质疑现行做法，例如，许多支付方要求所有初诊患者接受体格检查，这能否降低其未来几年患病的风险或医疗费用？

研究者往往有几个潜在的研究问题。同时参与不止一个项目有所帮助：一个项目停滞时，例如等待机构审查委员会（Institutional Review Board，IRB）或合作者回复时，可以推进另一个项目。但是，项目过多会分散有限的时间和资源。明智的办法是根据下文将讨论的优良问题特征，把最值得做的几个排出优先次序。有时，具有时限的机会，例如资助申请截止日，会要求你优先关注某个问题。可以把其他想法列下来，等精力允许时再回头处理。把问题与其他研究者、专科培训医师或学生分享，也可能使你以共同研究者或导师身份参与研究。

\subsection{专业知识}
\sourcepage{45}
成为研究领域的专家十分重要。扎实的学术积累是成功研究的基础。初学者应全面检索与研究问题有关的已发表文献，并批判性地阅读重要论文。下一步可以开展\term{系统综述}（systematic review），以建立和巩固对某一领域的专门知识；其文献综述也可作为资助申请及研究报告的背景材料。

在某领域积极开展工作的研究者，往往早在最新进展发表前就已知晓。掌握一个主题，需要持续了解多种信息来源，包括通过 PubMed 和期刊获取相关文献提醒、参加学术会议、与领域专家建立联系，甚至在社交媒体平台（例如 Twitter）关注有影响力的学者。可以查阅 ClinicalTrials.gov，了解是否已有与你设想相似的研究方案注册。该网站虽然主要涉及临床试验，有时也包括观察性研究。同样，如果打算进行系统综述，可查阅 Prospero 网站；许多系统综述已在该网站注册（\url{https://www.crd.york.ac.uk/prospero/#aboutpage}）。

\subsection{回应医疗或公共卫生需求}
有影响力的研究成果可能源于对更好治疗、检查或卫生服务的需求。对相关疾病、人群及项目的亲身经验，最有助于发现能够切实解决问题的机会。这往往也是开展新颖且有影响力研究的机会。

\subsection{留意新想法与新技术}
参加展示最新研究的会议很有帮助。与正式报告至少同样重要的是，海报交流和会间休息提供了与其他科学家非正式交谈的机会。研究者若能克服顾虑，在喝咖啡的间隙主动与报告者交流，可能收获颇丰，有时还会结识一位新的资深同行。如果某位报告者的工作与你特别相关，不妨先阅读她近期发表的论文，再与她联系，约好会议期间见面。

\sourcepage{46}
对普遍接受的观念保持怀疑，可以激发好的研究问题。例如，人们曾普遍认为，裂伤一旦穿透真皮，就必须缝合，才能迅速愈合并取得满意的外观。然而，Quinn 等根据亲身经验和病例系列资料注意到，中等大小的伤口无论创缘是否对合，都能够自行修复（2）。他们开展了一项随机试验，对长度小于 2 cm 的无并发症手部裂伤患者，用自来水冲洗伤口，并敷以抗生素敷料 48 小时。一组随机接受缝合，另一组不缝合。缝合组报告急诊治疗更疼痛、更耗时，但采用盲法评定后，两组的愈合时间及外观效果相似。

应用新技术，常会使人对熟悉的临床问题产生新的见解和疑问，进而形成新的\term{范式}（paradigms）（3）。例如，遗传学、分子技术、影像及数字健康技术的进展，催生了许多转化研究，并带来改变临床医学的治疗和检查。一个例子是使用\term{氘标记肌酸}（deuterated creatine，D3Cr）测量肌肉量；这种生物标志物比双能 X 线吸收测定扫描更准确，也更便宜（4）。D3Cr 测量可用于研究肌肉量能否预测择期手术后的结局，也可更准确地研究运动是否增加骨骼肌量。把某一领域的新概念、技术或发现用于另一领域的问题，能够产生创新的研究问题。

临床经验为探寻疾病的潜在原因和治疗方式提供了丰富的问题来源。仔细观察患者，可以识别新的危险因素、罕见遗传综合征或治疗并发症。将发现详细写成\term{病例报告}（case report）或\term{病例系列}（case series）发表，可以提醒其他人识别更多病例。
\sourcepage{47}
例如，临床医师在九名长期服用骨质疏松治疗药物阿仑膦酸的患者中，观察到不常见的股骨骨折（5）。这一发现最终促成大规模流行病学研究，证实阿仑膦酸使用时间越长，这类非典型骨折的风险越高（6）。

与其他研究者合作，常会产生可以进一步探究的新想法和新问题；当合作者带来另一领域的知识或方法时，尤其如此。这样的合作者作为团队成员，可以持续贡献创新和想法。

教学也是极好的灵感来源。准备报告，或者与善于提问的学生讨论时，往往会产生研究设想。学生常会要求解释常规做法或观念，从而暴露其中的不确定性或依据不足，指出需要研究之处。

构思研究问题、设想回答老问题的新方法以及自由探索想法，都离不开创造力。有些创意来自午餐时与同事的闲谈；另一些来自小组中对最新研究或自己想法的讨论。许多灵感则在独处时出现，例如准备讲座、淋浴、浏览社交媒体，或者静坐思考时。诀窍是让一个尚未解决的问题清晰地呈现在眼前，再让思绪围绕它自由展开。有意识地接触不同学科，例如阅读其他学科的期刊论文、聆听自己专业兴趣之外的会议报告，有时会产生跨领域的新想法，为自己的研究带来新的测量方法或概念。

害怕批评或担心想法与众不同，会抑制创造力，妨碍探究那些挑战现有范式的新问题。应寻求喜欢设想新点子的同事支持。同时还需要坚韧不拔，反复回到棘手的问题上，直到找到解决办法。

\section{好的研究问题具有什么特征}
能够形成良好研究计划的研究问题，应当可行（Feasible）、重要（Important）、新颖（Novel）且合乎伦理（Ethical）；
\sourcepage{48}
其首字母组成便于记忆的 FINE（表 2.1）。
\begin{bookbox}{表 2.1\quad 好的研究问题及研究计划的 FINE 标准}
\textbf{可行}
\begin{itemize}\item 参与者人数充足\item 技术能力充分\item 时间和经费负担得起\item 能获得资助\end{itemize}
\textbf{重要}
\begin{itemize}\item 可能改善临床医疗或公共卫生\item 有助于研究者因缜密而有价值的研究建立声誉\end{itemize}
\textbf{新颖}
\begin{itemize}\item 提供新发现\item 证实、反驳或扩展以往发现\item 可能在健康与疾病的概念、医疗实践或研究方法方面带来创新\end{itemize}
\textbf{合乎伦理}
\begin{itemize}\item 符合机构审查委员会的标准\item 关注健康公平\end{itemize}
\end{bookbox}

\subsection{可行}
最好及早了解研究某个问题的实际限制和困难，以免在不可行的方向上浪费大量时间和精力。
\begin{description}[leftmargin=2em,labelindent=0pt,style=nextline]
\item[参与者人数。]许多研究未能实现预期目的，是因为无法招募足够的参与者。应尽早估计所需样本量（第 5、6 章），考虑可能获得多少参与者，有多少人会被排除或拒绝入组，以及多少人可能失访。即使计划周密，估计也常常过于乐观。
\sourcepage{49}
在投入精力完善其他细节之前，应先确认有足够多符合条件且愿意参加的人。有时需要开展小规模先导调查，或检查已有的电子数据，才能确定。如果人数看来不足，可以考虑扩大纳入标准、取消不必要的排除标准、延长招募时间、增加参与者来源、采用更精密的测量方法、邀请同行参与多中心研究，或者改用另一种研究设计。
\item[必需的专业能力。]研究者必须具有设计研究、招募参与者、测量变量以及管理和分析数据所需的技能、设备与经验。与专家合作，或组建具有相应技能和经验的团队，能够补足研究者不熟悉的技术。比如，通常应从研究计划开始，就让统计学家加入团队，为样本量和分析方法提供建议。需要新测量方法时，例如检测一种新的生物标志物，应寻求熟悉最佳方法的专家。
\item[合适数据是否可用。]如果回答问题需要使用二手数据集（第 16 章），就应了解可能相关的数据集有哪些可用信息，以及这些信息的局限。
\item[时间与经费成本。]应估算项目每个部分的成本，并记住实际所需时间和经费通常会超过最初预计。如果预计成本超出现有经费，只能考虑更便宜的设计，或寻找额外资助。及早认清研究过于昂贵或耗时，可以在投入大量精力之前修改或放弃计划。
\item[研究范围。]研究者企图完成太多任务时，往往会出现问题，例如，为回答过多问题，反复接触一大批参与者并进行许多测量。
\sourcepage{50}
解决办法是缩小研究范围，只关注最重要的目标。许多科学家难以放弃回答有趣附带问题的机会，但这样做可能换来对主要问题更好的回答。
\item[获得资助的可能性。]很少有研究者能完全依靠个人或机构资源资助自己的研究，尤其是在需要招募和随访参与者、或进行昂贵测量时。如果无人承担实施费用，再精巧的研究设计也无法实现。第 20 章讨论如何寻找可能的资助来源。
\end{description}

\subsection{重要}
人们选择某个研究问题的动机很多：探明事实真相的兴奋、职业发展的顺理成章的下一步，或者获得支持以延续自己的研究方向。
\begin{description}[leftmargin=2em,labelindent=0pt,style=nextline]
\item[意义。]追究研究问题最好的动机，是改善健康与福祉。因此，应设想研究可能出现的各种结果，考虑每一种结果将如何增进科学知识、影响临床实践或公共卫生政策，或指导后续研究。还要确认，认为这个问题重要的不只有你自己。在投入大量精力拟订研究计划或资助申请之前，应先与导师、外部专家及潜在资助方的代表，例如美国国立卫生研究院（National Institutes of Health，NIH）的项目官员交谈，以免同行及资助机构认为这项研究不会改变实践。“意义”是资助申请单独评分的一个项目，而且常常决定项目的总体影响评分。
\item[促进职业发展。]另一个重要考虑是，探究这个研究问题是否有利于你的职业发展。研究能否培养新技能、建立合作关系，或者带来经费以扩大团队？结果能否让他人认可你研究的价值或缜密程度？相反，一个问题如果需要多年才能产生结果，而预期影响或关注度很低，就会拖慢职业发展。
\sourcepage{51}
仅仅为了在简历或晋升材料里增加论文而做研究，虽然有一定诱惑，最终却可能只积累了一批对科学和健康影响很小的成果。
\end{description}

\subsection{新颖}
一项研究如果只是重复已经确立的认识，就不值得投入精力和经费，也不大可能获得资助。判断拟议研究是否新颖，可以查阅文献，咨询了解尚未发表研究的专家，并在 NIH 的 Research Portfolio Online Reporting Tools（RePORT）网站（\url{http://report.nih.gov/categorical_spending.aspx}）及 ClinicalTrials.gov 网站搜索兴趣领域中已获资助项目的摘要。NIH 在评审研究申请时，重视其是否具有创新性：如果研究成功，能否通过新概念、新方法或新干预，改变研究或临床实践的范式（第 20 章）。

新颖固然重要，但研究问题不必完全原创。检验既往观察能否重复、一个人群中的发现是否也适用于其他人群，或新的测量方法能否澄清已知危险因素与疾病之间的关系，都可能值得研究。如果确证性研究能够避免以往研究的缺点，或者先前结果出人意料，其价值尤其大。

\subsection{合乎伦理}
合乎伦理的研究，不只是避免不当行为或伤害，还包括促进公平以及参与者和公众的福祉。
\begin{description}[leftmargin=2em,labelindent=0pt,style=nextline]
\item[涉及人的研究。]好的研究问题必须能够用合乎伦理的研究来回答。如果研究带来不可接受的风险，尤其是对易受伤害群体，或侵犯隐私（第 7 章），研究者就必须寻找其他回答方法。如果不确定研究或其中某个方面是否合乎伦理，应及早与 IRB 的代表讨论。
\sourcepage{52}
\item[健康公平。]制度性及其他结构性障碍是健康差异的重要原因。研究不应扩大不公平，而应针对并减轻这些差异。改善健康公平的一步，是额外努力纳入具有不同背景和经历的参与者，例如不同性别、种族或族裔、社会经济地位、英语熟练程度的人。另一步，是让传统上研究不足的人群参与研究设计与成果传播。第 15 章讨论如何通过与社区成员合作的研究框架，消除系统性不公平。
\end{description}

\section{制订研究计划}
尽早用一句话概括研究问题和研究设计，很有帮助。有些研究是描述性的，例如“对得克萨斯州幼儿园儿童花生过敏患病率的横断面研究”；但大多数研究旨在确定某个预测因素是否与结局有关，或是否造成结局，例如“估计光疗对北加利福尼亚足月新生儿两月龄时母乳喂养影响的队列研究”。

下一步是写一份简短的一页研究计划提纲，说明研究的重要性、设计、研究参与者及其选择方式、所需人数，以及拟作的测量（见第 1 章附录 1A）。这需要一点自律，却能迫使研究者澄清想法，找出需要处理的具体问题。提纲也为导师和同事提出具体建议提供依据。

制订研究计划是反复改进研究设计的过程，包括调整选择标准或重新估算样本量、与同行讨论改动，以及在适当时对关键环节进行预试验。随着方案成形，为确认有足够多符合条件且愿意参加的人而开展的先导研究，可能促使研究者修改招募计划。
\sourcepage{53}
与此同时，她也可能正在组建团队，以协助测量、制订统计计划、撰写资助申请（第 20 章），并在一切顺利时实施研究。

\subsection{主要问题与次要问题}
许多研究不止一个研究问题。临床试验常会考察干预对多个结局的作用。例如，妇女健康倡议（Women’s Health Initiative）的设计目标，是确定减少膳食脂肪摄入是否降低乳腺癌风险；一项重要的次要假设则是考察其对冠状动脉事件的作用（7）。几乎所有队列和病例对照研究都会针对每种结局考察几个危险因素。在设计中纳入多个研究问题，优点是可以提高效率，一项研究得到多个答案；缺点是增加设计、实施以及多重假设下统计推断的复杂性（第 5 章）。合理的策略是先确定一个主要研究问题，围绕它制订计划和估算样本量。再加入次要研究问题、测量或生物样本，使研究能够回答不止一个问题，并建立数据库和样本库，供进一步分析、撰写论文及开展支持后续研究申请的前期工作。

\section{选择导师并与导师合作}
构思研究问题和完善研究计划，需要作出许多判断，经验在其中不可替代。因此，初学者的一项必要策略，是跟随一位有经验、愿意且能够定期投入时间指导她的主要导师学习。通常还需要不同学科的共同导师组成团队，帮助她掌握研究内容及研究方法，获得开展项目的技能，并成长为独立研究者。

选择导师应慎重。可以与她以前指导过的人交谈，查看候选导师的简历，了解她是否曾与学生共同发表论文或取得资助，并与她面谈，了解她喜欢的工作方式，以及对受指导者的期望。
\sourcepage{54}
好的导师能够定期会面，并随时就你的想法进行非正式讨论。导师可能为你打开建立联系和争取资助的机会，鼓励独立工作，并在适当时让初学者参与资助申请和论文项目。这类良好关系还可能带来具体资源，例如研究助理或行政助理的协助；接触临床人群、数据集和样本库的机会；专业实验室及办公场所；统计支持；经费和研究团队，更不用说终生的友谊。

相反，不好的导师可能以压制创造力的方式批评新想法，从而妨碍成功；也可能把初学者工作的发现据为己有，或在发表、报告她的成果时自行担任主角。导师关系也受到性别、种族或族裔及权力所带来的社会与人际因素影响。如果年轻研究者觉得关系中存在虐待，应向其他资深同事求助；若担心骚扰、歧视或报复，则可向所在机构负责美国《教育法修正案》第九条（Title IX）的官员求助。更常见的是，许多导师只是太忙或注意力分散，无法顾及初学者的需求。无论哪种情形，如果与导师讨论仍无结果，我们建议寻找转向更合适指导者的办法，必要时请一位中立的资深同事或导师团队其他成员介入。

有时，导师因为缺少时间或兴趣，会交给你一个尚未完成的项目或稿件，往往来自以前的学生。这可能是以主要作者身份发表论文的好机会，也可能只是令人沮丧、对职业或专长毫无帮助的任务。（上一位学生放弃项目，也许自有原因。）初学者即使觉得难以启齿，也应讨论它对自己训练和职业的潜在价值；
\sourcepage{55}
如果收益很小，就应谢绝。这种情形下，导师团队可以对值得怀疑的任务或项目起到制约作用。

\section{小结}
所有研究都应从一个研究问题开始，明确研究者希望知道什么；目标是找到一个能够发展成良好研究计划的问题。
\begin{enumerate}
\item 专业知识是提出值得研究的问题的必要条件。系统综述与兴趣领域有关的研究，是很好的起点。参加会议、结识其他科学家，并在社交媒体关注感兴趣的领域，可以让知识范围超出已发表文献。
\item 好的问题来自与同行交谈、批判性思考临床做法及问题、用新方法处理旧问题，以及从教学、自由遐想和坚持不懈解决棘手问题的过程中获得的想法。
\item 在投入大量时间撰写申请或开展研究之前，应判断问题和计划是否符合 FINE：可行、重要、新颖且合乎伦理。资助方会优先考虑可能对科学和健康产生创新性及重要影响的研究。
\item 应及早把研究问题写成一页提纲，具体说明研究重要性、设计、参与者、所需人数和拟作测量。
\item 发展研究问题和研究计划是反复改进的过程，包括请教导师与同行，并组建协助测量、统计及资助申请的团队。
\item 大多数研究有多个问题。设计和实施时，关注一个主要问题很有帮助。添加次要问题和测量，并储存生物样本，可以为后续论文及研究问题的前期研究积累资源。
\sourcepage{56}
\item 初学者必须选择一两位资深科学家作为导师：他们应有经验，愿意会面、提供资源和联系、鼓励创造力，并促进年轻科学家的独立与学术能见度。
\end{enumerate}

\begin{references}
\item Perissinotto CM, Cenzer IS, Covinsky K. Loneliness in older persons: a predictor of functional decline and death. Arch Intern Med. 2012;172(14):1078-1083.
\item Quinn J, Cummings S, Callaham M, et al. Suturing versus conservative management of lacerations of the hand: randomized controlled trial. BMJ. 2002;325:299-301.
\item Kuhn TS. The Structure of Scientific Revolutions. University of Chicago Press; 1962.
\item Cawthon P, Orwoll ES, Peters KE, et al. Strong relation between muscle mass determined by D3-creatine dilution, physical performance, and incidence of falls and mobility limitations in a prospective cohort of older men. J Gerontol A Biol Sci Med Sci. 2019;74(6):844-852.
\item Goh SK, Yang KY, Koh JS, et al. Subtrochanteric insufficiency fractures in patients on alendronate therapy: a caution. J Bone Joint Surg Br. 2007;89(3):349-353.
\item Black DM, Geiger EJ, Eastell R, et al. Atypical femur fracture risk versus fragility fracture prevention with bisphosphonates. N Engl J Med. 2020;383(8):743-753.
\item Howard BV, Van Horn L, Hsia J, et al. Low-fat dietary pattern and risk of cardiovascular disease: the Women’s Health Initiative Randomized Controlled Dietary Modification Trial. JAMA. 2006;295(6):655-666.
\end{references}

\sourcepage{57}
\chapterappendix{2A}{第 2 章习题：构思研究问题与制订研究计划}
\begin{enumerate}
\item 考虑研究问题：“大麻使用与健康有什么关系？”首先，将其改写为更具体的表述，明确研究设计、预测变量、结局变量和目标人群。随后讨论你选择的问题与设计是否符合 FINE 标准（可行、重要、新颖、合乎伦理）。改写问题和设计，解决不符合这些标准之处。
\item 考虑研究问题：“对乙酰氨基酚（acetaminophen，又称扑热息痛，paracetamol）是否引起哮喘？”设想自己回到 2000 年，即这个问题开始受到关注的时候，分别用一句话描述回答该问题的两项观察性研究及一项临床试验。每项描述都须明确设计、预测变量、结局变量和目标人群。然后分别考虑问题与设计是否符合 FINE 标准。
\item 运用本章观点及自己的兴趣，构思一个研究问题，并为自己可能开展的研究写一页提纲。它是否符合 FINE 标准？讨论不同的设计、目标人群、样本和变量，以改进研究。
\end{enumerate}
